\clearpage
\fancyhead[L]{\small JUNO. ENTDECKUNG.}
\hypertarget{native-258-lower}{}
\begin{center}
EINIGE NACHRICHTEN ÜBER DEN NEUEN PLANETEN.

\rule{0.75\linewidth}{0.4pt}

Braunschweigisches Magazin. Stück 40. 1804, 6. Oktober.

\rule{0.75\linewidth}{0.4pt}
\end{center}

Die kürzlich in \emph{Lilienthal} gemachte Entdeckung eines neuen beweglichen Sterns, über die bereits einige Zeitungsberichte erschienen sind, hat die öffentliche Neugier so sehr geweckt, dass Menschen, denen der Fortschritt der Wissenschaften am Herzen liegt, einige genauere Nachrichten über diese bemerkenswerte Entdeckung wahrscheinlich begrüßen werden. Das gilt umso mehr, als sich
\SourceLeaf{259}
jetzt, nachdem der Lauf des Sterns seit drei Wochen beobachtet wird, schon recht zuverlässige Schlüsse über seine wahre Bahn ziehen lassen. Die meisten Leser dieser Blätter würden diese Ergebnisse sonst erst viel später zu Gesicht bekommen.

Herr \textsc{Harding}, Mitarbeiter des berühmten Justizrats \textsc{Schröter} in Lilienthal, arbeitete bereits seit einem halben Jahr an einem astronomischen Atlas eigener Art. Er soll nur die Teile des Fixsternhimmels umfassen, in denen die beiden neu entdeckten Planeten \emph{Ceres} und \emph{Pallas} sichtbar sein können. Zum Auffinden und Beobachten dieser kleinen Himmelskörper sind sehr detaillierte Sternkarten unbedingt nötig, die wenigstens bis zur achten Größe hinab alle Sterne vollständig enthalten. Sämtliche bisher vorhandenen Himmelskarten reichen dazu bei Weitem nicht aus. Das Vorhaben von Herrn \textsc{Harding} verdient daher große Anerkennung. Von der großen Sorgfalt und dem unermüdlichen Fleiß, mit denen er es ausführt, zeugen einige ausgezeichnete Detailkarten, die er bereits für Herrn von \textsc{Zach}s Monatliche Correspondenz geliefert hat. Herr \textsc{Harding} stützt sich dabei hauptsächlich auf die Histoire céleste française von \textsc{Lalande}, ein unschätzbares Werk. Es enthält die innerhalb von 10 Jahren an der Sternwarte der Militärschule in Paris hauptsächlich von \textsc{Lalande}s Neffen \textsc{Michel Lefrançais Lalande} gemachten Beobachtungen von 50000 Fixsternen. Alle Sterne daraus, die in die von Herrn \textsc{Harding} bearbeiteten Zonen fallen, trägt er in seine Karten ein. Jeden einzelnen sucht er aber vorher selbst am Himmel auf, um sich zu vergewissern, dass er wirklich vorhanden ist, und um etwaige Beobachtungs-, Schreib- und Druckfehler zu berichtigen, die bei einer so gewaltigen Arbeit unvermeidlich sind. Damit begnügt sich Herr \textsc{Harding} noch nicht. Wo er weitere lichtschwächere Fixsterne bis wenigstens zur achten Größe findet, die in der Histoire céleste fehlen, bestimmt er sie selbst und trägt sie in seinen Karten nach. Solche Sterne gibt es noch in großer Zahl.

Gerade bei einer solchen Durchmusterung einer Himmelsgegend stieß Herr \textsc{Harding} am 1. September im Sternbild der Fische auf einen Stern achter Größe, der in der Histoire céleste fehlt. Weil das unzählige Male vorkommt, schöpfte er keinen Verdacht und trug ihn vorläufig allein nach Augenmaß als Fixstern in seine Karten ein. Als er jedoch drei Tage später wieder nachsah, stand an jener Stelle kein Stern mehr. In der Nähe, etwa einen halben Grad südwestlich davon, zeigte sich dagegen ein ganz ähnlicher Stern an einer Stelle, an der er sich nicht erinnerte, am 1. September einen Stern gesehen zu haben. Das erregte starken Verdacht, ließ aber noch keinen sicheren Schluss zu, weil es vielleicht nur an einer bei lichtschwachen Fixsternen sehr gewöhnlichen Helligkeitsänderung oder an einem Irrtum lag. Am 5. September schließlich waren beide bisherigen Stellen leer, während an einer Stelle in entsprechendem Abstand nach Südwesten ein Stern stand. Herr \textsc{Harding} konnte nun nicht mehr daran zweifeln, dass die an den drei Abenden wahrgenommenen Sterne ein und derselbe waren. Er nahm daher eine regelrechte Beobachtung des entdeckten beweglichen Sterns vor und setzte seine Beobachtungen an den folgenden Abenden fort. Sie bestätigten eine ganz regelmäßige rückläufige Bewegung, die den Stern zugleich nach Süden führte. Er unterrichtete auswärtige Astronomen von seiner wichtigen Entdeckung.

Ich erhielt die Nachricht am 11., als der Himmel hier am Abend ganz bedeckt war. Am 12. dagegen war es sehr klar. Ich fand \textsc{Harding}s Stern sofort und bestimmte seinen Ort durch eine gute Beobachtung. Seitdem konnte ich ihn bis heute (30. Sept.) noch an elf verschiedenen Abenden beobachten, zuletzt am 28. Seine Bewegung ist noch immer rückläufig und nach Süden gerichtet. Der bisherige Bogen kommt einer geraden Linie so nahe, dass die Krümmung kaum zu bemerken ist, doch die Geschwindigkeit hat merklich zugenommen. Diese Erscheinungen vertragen sich sehr gut mit der Annahme, dass dieser Himmelskörper wirklich eine Planetenbahn beschreibt. So allgemein betrachtet, lassen sie natürlich zunächst nur Vermutungen zu. Gewissheit kann uns nur eine
\SourceLeaf{260}
Rechnung geben, die mit so viel Genauigkeit und Vorsicht durchgeführt wird, wie eine so heikle und schwierige Untersuchung erfordert.

Im Aussehen ähnelt dieser Stern Ceres und Pallas sehr, ist aber noch etwas heller als sie. Nach meiner Schätzung entspricht seine Helligkeit ungefähr der eines Fixsterns $7^{ter}$ bis $8^{ter}$ Größe. Selbst am 18. September, als der fast volle Mond nicht weit von ihm entfernt stand, konnte ich ihn mit einem Spiegelteleskop von \textsc{Short} ohne die geringste Anstrengung sehen und beobachten. Bei diesem Stern zeigt sich nicht die geringste Spur eines Nebels, wie er bei Kometen vorkommt, auch mit den großen Lilienthaler Instrumenten nicht.

Von den anderen auswärtigen Astronomen, die diesen Stern bereits beobachten, kenne ich bisher nur Herrn Dr. \textsc{Olbers}, Herrn von \textsc{Zach} und Herrn Prof. \textsc{Bode}. Der erste begann mit den Beobachtungen am 7., Herr von \textsc{Zach} am 13. und Herr \textsc{Bode} am 21. Da der Stern aber noch mindestens bis Januar 1805 sichtbar bleiben wird, hoffe ich zuversichtlich, dass er noch an allen europäischen Sternwarten beobachtet werden kann. So kurze Zeit dieser bemerkenswerte Himmelskörper auch bisher beobachtet worden ist, habe ich doch schon versucht, vorwiegend anhand meiner eigenen Beobachtungen eine erste annähernde Bestimmung der Hauptdimensionen seiner wahren Bahn vorzunehmen. Das befriedigt zum Teil und fürs Erste die ganz natürliche allgemeine Wissbegier. Einen wesentlicheren Nutzen hat eine solche Arbeit aber dadurch, dass sie uns in die Lage versetzt, den Lauf des Sterns schon für einige Zeit \emph{voraus}zusagen. So können wir weiter entfernten Astronomen, die die Nachrichten über den bereits zurückgelegten Weg womöglich erst so spät erhalten, dass sie den Stern damit nicht mehr unmittelbar auffinden können, ungefähr zeigen, wo sie ihn suchen müssen, und ihnen vielleicht eine vergebliche Suche ersparen.

Das bemerkenswerte Ergebnis dieser Untersuchung war, dass \textsc{Harding}s Stern ebenso wie Ceres und Pallas eine Planetenbahn zwischen Mars und Jupiter beschreibt. Die Umlaufzeit und die mittlere Entfernung von der Sonne sind fast (oder ganz) genau dieselben wie bei diesen beiden. Die Exzentrizität ist noch etwas größer als bei Pallas, die Bahnneigung dagegen viel kleiner als bei Pallas und nur wenig größer als bei Ceres. Es ergaben sich eine Umlaufzeit von 1662 Tagen, eine mittlere Entfernung von 2.746, eine Exzentrizität von 0.287, eine Bahnneigung von 13 Grad 35 Minuten, eine Länge des Aphels von 239 Grad 14 Min., eine Länge des aufsteigenden Knotens von 171 Grad 16 Min. Die mittlere Länge in der Bahn betrug am 5. Sept. zu Mittag im Meridian der Sternwarte Seeberg 20 Grad 39 Min.

Nach diesen Ergebnissen kann es jetzt schon als ausgemacht gelten, dass \textsc{Harding}s Stern ebenso wie die anderen Planeten seine regelmäßigen Zeiten der Sichtbarkeit und Unsichtbarkeit haben und von nun an ein ständiger Gegenstand unserer Beobachtungen sein wird. Ihm gebührt daher mit vollem Recht der Rang eines Planeten. Herr \textsc{Harding} hat deshalb bereits den Namen \emph{Juno} für ihn vorgeschlagen, um ihn den älteren Planeten gleichzustellen. Zweifellos werden alle Astronomen das Recht des Entdeckers achten und seiner Namenswahl zustimmen. Als Zeichen der Juno ist ihr Zepter mit einem Stern vorgeschlagen worden.

Juno steht jetzt innerhalb des gut erkennbaren Vierecks, das die vier Sterne in den Fischen bilden, die in \textsc{Flamsteeds} Verzeichnis mit 27, 29, 30, 33 bezeichnet sind. Wer gute Augen hat, sieht diese Sterne bei klarem Himmel ohne Fernrohr. Juno wird ihre Bewegung noch einige Zeit in der bisherigen Richtung fortsetzen, aber mit schon jetzt wieder abnehmender Geschwindigkeit. In der zweiten Oktoberhälfte wird die Krümmung des Wegs nach Südosten immer deutlicher. Etwa am 30. Oktober geht die rückläufige Bewegung in eine rechtläufige über, und etwa am 17. November, wenn Juno eine südliche Deklination von $11\frac{1}{4}$ Grad erreicht haben wird, beginnt diese wieder abzunehmen, und der Weg biegt wieder nach Norden ab.

\SourceLeaf{261}
Vielleicht gibt es unter den Lesern dieser Blätter den einen oder anderen, der über geeignete Instrumente und genügend Geschick verfügt, um zu versuchen, diesen neuen Planeten aufzuspüren und ihn auf seinem Lauf zu verfolgen. Deshalb füge ich noch die genauen Angaben sowohl zu seinen bisherigen als auch zu seinen künftigen Positionen hinzu, jene nur von Woche zu Woche, diese von 3 zu 3 Tagen. Alle Angaben gelten für Mitternacht. [Diese Ephemeride ist ein Auszug aus der oben auf S. 253 abgedruckten.]
